\providecommand{\Res}[2]{\qty[round-mode=figures,round-precision=3]{\fpeval{#1}}{#2}}
\providecommand{\Val}[1]{\num[round-mode=figures,round-precision=4]{\fpeval{#1}}}
\providecommand{\ResSci}[2]{\qty[round-mode=figures,round-precision=3,exponent-mode=scientific]{\fpeval{#1}}{#2}}
Pour de petits angles, $\sin\theta \approx y/L = \Res{\vsi{y}/\v{L}}{}$.

\subsubsection*{Question a) Facteur d'interférence}
\begin{enumerate}[label=\arabic*)]
    \item \textbf{Variables connues et recherchées :}
    \begin{itemize}[label= -]
        \item $\lambda = \qty{\vsi{lam}}{m}$, $d = \qty{\vsi{sep}}{m}$, $\sin\theta = \Val{\vsi{y}/\v{L}}$
        \item Variable recherchée : $\cos^2(\phi/2)$
    \end{itemize}
    \item \textbf{Équations pertinentes :}
    \[ \frac{\phi}{2} = \frac{\pi d\sin\theta}{\lambda} \]
    \item \textbf{Isolation de la variable recherchée :} on calcule $\phi/2$, puis $\cos^2(\phi/2)$ (calculatrice en radians).
    \item \textbf{Résolution numérique :}
    \[ \frac{\phi}{2} = \frac{\pi \cdot \num{\vsi{sep}} \cdot \Val{\vsi{y}/\v{L}}}{\num{\vsi{lam}}} = \Res{pi*\vsi{sep}*\vsi{y}/(\vsi{lam}*\v{L})}{rad} \quad\Rightarrow\quad \cos^2(\phi/2) = \Res{cos(pi*\vsi{sep}*\vsi{y}/(\vsi{lam}*\v{L}))^2}{} \]
    \item \textbf{Réponse :} le facteur d'interférence vaut \textbf{\Res{cos(pi*\vsi{sep}*\vsi{y}/(\vsi{lam}*\v{L}))^2}{}}.
\end{enumerate}

\subsubsection*{Question b) Facteur de diffraction}
\begin{enumerate}[label=\arabic*)]
    \item \textbf{Variables connues et recherchées :} $a = \qty{\vsi{larg}}{m}$ ; variable recherchée : $(\sin\alpha/\alpha)^2$
    \item \textbf{Équations pertinentes :}
    \[ \alpha = \frac{\pi a\sin\theta}{\lambda} \]
    \item \textbf{Isolation de la variable recherchée :} on calcule $\alpha$, puis $(\sin\alpha/\alpha)^2$.
    \item \textbf{Résolution numérique :}
    \[ \alpha = \frac{\pi \cdot \num{\vsi{larg}} \cdot \Val{\vsi{y}/\v{L}}}{\num{\vsi{lam}}} = \Res{pi*\vsi{larg}*\vsi{y}/(\vsi{lam}*\v{L})}{rad} \quad\Rightarrow\quad \left(\frac{\sin\alpha}{\alpha}\right)^2 = \Res{(sin(pi*\vsi{larg}*\vsi{y}/(\vsi{lam}*\v{L}))/(pi*\vsi{larg}*\vsi{y}/(\vsi{lam}*\v{L})))^2}{} \]
    \item \textbf{Réponse :} le facteur de diffraction vaut \textbf{\Res{(sin(pi*\vsi{larg}*\vsi{y}/(\vsi{lam}*\v{L}))/(pi*\vsi{larg}*\vsi{y}/(\vsi{lam}*\v{L})))^2}{}}.
\end{enumerate}

\subsubsection*{Question c) Intensité au point P}
\[ I = I_0\cos^2\left(\frac{\phi}{2}\right)\left(\frac{\sin\alpha}{\alpha}\right)^2 = \num{\v{I0}} \cdot \Val{cos(pi*\vsi{sep}*\vsi{y}/(\vsi{lam}*\v{L}))^2} \cdot \Val{(sin(pi*\vsi{larg}*\vsi{y}/(\vsi{lam}*\v{L}))/(pi*\vsi{larg}*\vsi{y}/(\vsi{lam}*\v{L})))^2} = \Res{\v{I0}*cos(pi*\vsi{sep}*\vsi{y}/(\vsi{lam}*\v{L}))^2*(sin(pi*\vsi{larg}*\vsi{y}/(\vsi{lam}*\v{L}))/(pi*\vsi{larg}*\vsi{y}/(\vsi{lam}*\v{L})))^2}{W/m^2} \]
L'intensité en P est \textbf{\Res{\v{I0}*cos(pi*\vsi{sep}*\vsi{y}/(\vsi{lam}*\v{L}))^2*(sin(pi*\vsi{larg}*\vsi{y}/(\vsi{lam}*\v{L}))/(pi*\vsi{larg}*\vsi{y}/(\vsi{lam}*\v{L})))^2}{W/m^2}}.
